% ── CHAPTER 12: CONVERGENCE 6: MOTION AS LIFE (~6,000 words) ──
\chapter{Motion as Life: Reality Is Fundamentally Dynamic}
\label{ch:motion-as-life}

% STATUS: TODO
% TARGET: ~6,000 words
% CONVERGENCE: 6 of 6 — Motion as Life
% PRIMARY SOURCES:
%   pub2/sections/05-six-convergences.tex Convergence6  — QM + process + Bahá'í
%   pub3/sections/03-six-convergences.tex Convergence6  — Bahá'í paragraph
% KEY CITATIONS: Heisenberg1927, Whitehead1929, PUP
% NOTE: "The most direct of the six" (pub3 §3.6). Two traditions arrive by
%   entirely different routes at: reality is dynamic at its deepest level;
%   inertness = nonexistence. Heisenberg uncertainty ↔ PUP Talk 52.
% HARD RULES: See BOOK_SHELL.md

% \epigraph{Creation is the expression of motion. Motion is life.}{\AbdulBaha,
%   \textit{The Promulgation of Universal Peace}}

%  OPENING (~500 words)
% TODO: Close with the most direct of the six convergences. The claim: there
% is no rest at the bottom of physical description. And the Bahá'í writings
% say the same thing from an entirely different direction.
% Transition note: this is the last of the six convergences before Part IV —
% which asks what paradigm is adequate to all six together.

\section{The Physics: Zero-Point Energy and the Restless Quantum}
\label{sec:ch12-physics}
% (~1,500 words)
% TODO: Accessible treatment of Heisenberg uncertainty and zero-point motion.
% KEY POINTS:
%   - Heisenberg uncertainty: perfectly localizing a particle's position would
%     require infinite momentum — infinite energy. Perfect rest is forbidden.
%   - Zero-point energy: even at absolute zero, a quantum system cannot be
%     at rest; it retains irreducible fluctuations
%   - The quantum vacuum is not empty; it seethes with persistent fluctuations
%   - Virtual particles: the vacuum is not a passive background but an active
%     medium of dynamic becoming
%   - There is no rest at the bottom of physical description
% ANALOGY: A spinning top — you can slow it down but you can never, in QM,
%   stop it completely.
% SOURCE: pub2 §5.6 QM paragraph; pub3 §3.6 physics paragraph
% CITE: Heisenberg1927

\section{The Process-Philosophical View: Process All the Way Down}
\label{sec:ch12-process}
% (~1,500 words)
% TODO: Whitehead on the primacy of process.
% KEY POINTS:
%   - Actual occasions are events of becoming, not enduring substances
%   - There are no truly static actual occasions — even the most "inert"
%     physical reality is a pattern of micro-events of becoming
%   - The apparent stability of macroscopic objects is real but derivative:
%     high-frequency societies of occasions maintaining a pattern
%   - Creativity (the ultimate metaphysical principle in Whitehead) is
%     perpetual novelty: each new occasion is genuinely new; nothing is
%     merely a repetition of what came before
% SOURCE: pub2 §5.6 process paragraph; pub2 §3.1 (primacy of process)
% CITE: Whitehead1929

\section{The Bahá'í View: Motion Is Life; Inertness Is Death}
\label{sec:ch12-bahai}
% (~1,500 words)
% TODO: Full development of PUP Talk 52.
% KEY QUOTE: "Creation is the expression of motion. Motion is life. A moving
%   object is a living object, whereas that which is motionless and inert is
%   as dead. All energy is contingent upon motion." ['Abdu'l-Bahá, PUP Talk 52]
% KEY POINTS:
%   - This is not a metaphor; it is an ontological claim
%   - Existence = process; inertness = nonexistence (or at minimum, near-death)
%   - The claim is universal: "all energy is contingent upon motion"
%   - This maps onto the Heisenberg result: at the quantum level, there is no
%     rest; all existence is perpetual dynamic process
%   - The structural isomorphism: BOTH traditions derive, from independent
%     starting points, the conclusion that reality is dynamic at its core
% SOURCE: pub3 §3.6; pub2 §5.6 Bahá'í paragraph; pub2 §4 (passage 2)
% CITE: PUP

\section{The Most Direct Convergence}
\label{sec:ch12-convergence}
% (~700 words)
% TODO: State the isomorphism. This is the most direct of the six: two
%   traditions arrive at the same conclusion (reality = motion; inertness =
%   non-being) by entirely different routes (Heisenberg commutation relations
%   vs. metaphysical reflection on the nature of creation).
%   End with transition to Part IV: the six convergences together demand a
%   new paradigm. That is what the next chapter asks for.

\begin{convergencemoment}{Motion as Life: Reality Is Fundamentally Dynamic}
% TODO: Fill in when drafting.
% From the quantum side: Heisenberg uncertainty forbids rest; zero-point
%   energy ensures perpetual motion at the bottom of physical description.
% From the process-philosophical side: actual occasions are events of
%   becoming; creativity ensures perpetual novelty; no static substance.
% From the Bahá'í side: "Creation is the expression of motion. Motion is
%   life... that which is motionless and inert is as dead." [PUP Talk 52]
% One structure — the most direct of the six: reality is dynamic at its
%   deepest level; inertness and nonexistence are the same thing.
\end{convergencemoment}
